% Options for packages loaded elsewhere
\PassOptionsToPackage{unicode}{hyperref}
\PassOptionsToPackage{hyphens}{url}
\documentclass[
  11pt,
]{article}
\usepackage{xcolor}
\usepackage[left=3cm,right=3cm,top=2.5cm,bottom=2.5cm]{geometry}
\usepackage{amsmath,amssymb}
\setcounter{secnumdepth}{5}
\usepackage{iftex}
\ifPDFTeX
  \usepackage[T1]{fontenc}
  \usepackage[utf8]{inputenc}
  \usepackage{textcomp} % provide euro and other symbols
\else % if luatex or xetex
  \usepackage{unicode-math} % this also loads fontspec
  \defaultfontfeatures{Scale=MatchLowercase}
  \defaultfontfeatures[\rmfamily]{Ligatures=TeX,Scale=1}
\fi
\usepackage{lmodern}
\ifPDFTeX\else
  % xetex/luatex font selection
\fi
% Use upquote if available, for straight quotes in verbatim environments
\IfFileExists{upquote.sty}{\usepackage{upquote}}{}
\IfFileExists{microtype.sty}{% use microtype if available
  \usepackage[]{microtype}
  \UseMicrotypeSet[protrusion]{basicmath} % disable protrusion for tt fonts
}{}
\usepackage{setspace}
\makeatletter
\@ifundefined{KOMAClassName}{% if non-KOMA class
  \IfFileExists{parskip.sty}{%
    \usepackage{parskip}
  }{% else
    \setlength{\parindent}{0pt}
    \setlength{\parskip}{6pt plus 2pt minus 1pt}}
}{% if KOMA class
  \KOMAoptions{parskip=half}}
\makeatother
\usepackage{longtable,booktabs,array}
\usepackage{caption}
\captionsetup[table]{skip=6pt}
\newcounter{none} % for unnumbered tables
\usepackage{calc} % for calculating minipage widths
% Correct order of tables after \paragraph or \subparagraph
\usepackage{etoolbox}
\makeatletter
\patchcmd\longtable{\par}{\if@noskipsec\mbox{}\fi\par}{}{}
\makeatother
% Allow footnotes in longtable head/foot
\IfFileExists{footnotehyper.sty}{\usepackage{footnotehyper}}{\usepackage{footnote}}
\makesavenoteenv{longtable}
\usepackage{graphicx}
\makeatletter
\newsavebox\pandoc@box
\newcommand*\pandocbounded[1]{% scales image to fit in text height/width
  \sbox\pandoc@box{#1}%
  \Gscale@div\@tempa{\textheight}{\dimexpr\ht\pandoc@box+\dp\pandoc@box\relax}%
  \Gscale@div\@tempb{\linewidth}{\wd\pandoc@box}%
  \ifdim\@tempb\p@<\@tempa\p@\let\@tempa\@tempb\fi% select the smaller of both
  \ifdim\@tempa\p@<\p@\scalebox{\@tempa}{\usebox\pandoc@box}%
  \else\usebox{\pandoc@box}%
  \fi%
}
% Set default figure placement to htbp
\def\fps@figure{htbp}
\makeatother
\setlength{\emergencystretch}{3em} % prevent overfull lines
\providecommand{\tightlist}{%
  \setlength{\itemsep}{0pt}\setlength{\parskip}{0pt}}
%% tools/stamp.tex
%% Provenance footer for PDFs rendered from this project.
%% Vendored by zzcollab; this file is not project-specific. The
%% footer prints the source document and the time of rendering.
%% The git version is defined here too, and so is available to a
%% project that wants it on the page, but it is left off the
%% default footer: it already names the staged copy in share/ and
%% appears in its manifest row. All values are supplied at render
%% time by tools/stamp-render.R, which writes a small generated
%% file that redefines the macros below. The \providecommand
%% defaults keep a bare LaTeX compile working when the document is
%% built without the wrapper.
\usepackage{fancyhdr}
\usepackage{xcolor}
\providecommand{\stampsource}{\detokenize{(source not recorded)}}
\providecommand{\stamptime}{(time not recorded)}
\providecommand{\stampversion}{\detokenize{(version not recorded)}}
\setlength{\footskip}{40pt}
\newcommand{\stampfootcontent}{%
  \footnotesize\color{gray}%
  \parbox[b]{\textwidth}{%
    \stamptime\hfill\thepage\\[2pt]%
    \ttfamily\raggedright\stampsource}}
\pagestyle{fancy}
\fancyhf{}
\fancyfoot[C]{\stampfootcontent}
\renewcommand{\headrulewidth}{0pt}
\renewcommand{\footrulewidth}{0.4pt}
%% The title page uses the `plain` page style; redefine it so the
%% stamp appears there too.
\fancypagestyle{plain}{%
  \fancyhf{}%
  \fancyfoot[C]{\stampfootcontent}%
  \renewcommand{\headrulewidth}{0pt}%
  \renewcommand{\footrulewidth}{0.4pt}}
\renewcommand{\stampsource}{\textasciitilde{}/\allowbreak{}prj/\allowbreak{}res/\allowbreak{}04-runin-power-analysis/\allowbreak{}runinpower/\allowbreak{}analysis/\allowbreak{}report/\allowbreak{}03-allocation/\allowbreak{}bullets.Rmd}
\renewcommand{\stamptime}{2026-08-15 16:32 PDT}
\renewcommand{\stampversion}{\detokenize{455c889-wip}}
\usepackage{amsmath}
\usepackage{amssymb}
\usepackage{booktabs}
\usepackage{longtable}
\usepackage{float}
\usepackage[mathlines]{lineno}
\linenumbers
\usepackage{bookmark}
\IfFileExists{xurl.sty}{\usepackage{xurl}}{} % add URL line breaks if available
\urlstyle{same}
% fallback for those not using the hyperref driver hyperxmp:
\makeatletter
\@ifundefined{xmpquote}{\newcommand{\xmpquote}[1]{#1}}{}
\makeatother
\hypersetup{
  pdftitle={Structured summary: \textbar{} Optimal Allocation of Observations across Run-In{,} Treatment{,} and Common Close Phases in Longitudinal Clinical Trials},
  pdfauthor={Ronald G. Thomas, Ph.D.},
  hidelinks,
  pdfcreator={LaTeX via pandoc}}

\title{Structured summary: \textbar{} Optimal Allocation of Observations across Run-In, Treatment, and Common Close Phases in Longitudinal Clinical Trials}
\author{Ronald G. Thomas, Ph.D.\thanks{Wertheim School of Public Health, UCSD; ORCID: 0000-0003-1686-4965}}
\date{August 15, 2026}

\begin{document}
\maketitle

{
\setcounter{tocdepth}{2}
\tableofcontents
}
\setstretch{1.1}
\emph{This document is a structured, section-by-section summary of the key points argued in the companion manuscript, ``| Optimal Allocation of Observations across Run-In, Treatment, and Common Close Phases in Longitudinal Clinical Trials.'' It is provided as a supplementary reading aid and does not stand on its own; all notation, derivations, simulation detail, and citations are in the main manuscript.}

\bigskip\noindent

\rule{\textwidth}{0.4pt}\bigskip

\section{Introduction}\label{introduction}

\textbf{Additional run-in and common close observations both improve efficiency but with diminishing or $J_1$-dependent returns.}

\begin{itemize}\setlength{\itemsep}{2pt}\setlength{\parskip}{0pt}
\item Run-in observations appended to a fixed treatment phase improve efficiency with diminishing returns.
\item Most of the attainable gain comes from the first two run-in observations; the marginal value falls quickly thereafter.
\item Common close gains depend on the number of treatment observations $J_1$.
\end{itemize}

\textbf{The efficiency results motivate an allocation question left open by Paper 1.}

\begin{itemize}\setlength{\itemsep}{2pt}\setlength{\parskip}{0pt}
\item Trials operate under constraints on total visits or total duration.
\item Paper 1 characterizes efficiency but not the optimal phase split.
\item The question is how to distribute observations across the three phases.
\end{itemize}

\textbf{Prior work addresses related design questions but not optimal allocation across run-in, treatment, and common close phases.}

\begin{itemize}\setlength{\itemsep}{2pt}\setlength{\parskip}{0pt}
\item Existing optima target time-point positioning or general design guidelines, not the run-in setting.
\item Multiple pre-treatment measurements are known to benefit ANCOVA-analyzed designs.
\item A substantial optimal-design literature treats measurement schedules for mixed models through the Fisher information matrix, but does not address a phase structure in which one phase carries its own nuisance parameter.
\item Available software computes slope-comparison power but does not optimize the phase allocation.
\end{itemize}

\subsection{Fixed total visits}\label{fixed-total-visits}

\textbf{The fixed-visit problem minimizes the treatment-effect variance over the phase split under a total-count constraint.}

\begin{itemize}\setlength{\itemsep}{2pt}\setlength{\parskip}{0pt}
\item The budget $J$ excludes baseline and must be split as $J_0 + J_1 + J_2 = J$.
\item Feasibility requires $J_0 \ge 0$, $J_1 \ge 1$, and $J_2 \ge 0$.
\item The objective is the $\gamma\gamma$ entry of $(X'\Sigma^{-1}X)^{-1}$.
\end{itemize}

\subsection{Fixed total duration}\label{fixed-total-duration}

\textbf{The fixed-duration problem instead constrains total trial length and may also optimize observation timing.}

\begin{itemize}\setlength{\itemsep}{2pt}\setlength{\parskip}{0pt}
\item Duration is fixed via $T = (J_0 + J_1 + J_2) \cdot t$, so the count constraint becomes $J_0 + J_1 + J_2 = T/t$.
\item The variance $\Var_1(\hat\gamma)$ is minimized subject to this constraint.
\item A general version permits non-equally-spaced observations and optimizes over the observation times.
\end{itemize}

\subsection{\texorpdfstring{The \(J_0=1\) anomaly}{The J\_0=1 anomaly}}\label{the-j_01-anomaly}

\textbf{In an intermediate signal-to-noise band the first run-in observation is less valuable than the second because it introduces a new model parameter.}

\begin{itemize}\setlength{\itemsep}{2pt}\setlength{\parskip}{0pt}
\item At $J_1 = 2$, $J_2 = 0$, $\sigma = t = 1$, $\Delta V_1 < \Delta V_2$ holds for $\sigma_b / \sigma$ between approximately 0.53 and 1.51, and fails outside this band.
\item Adding one run-in observation expands the model from two to three parameters, consuming a degree of freedom.
\item The second run-in observation adds no parameter and delivers pure information gain.
\end{itemize}

\textbf{The single-run-in penalty parallels a known ANCOVA result and yields a regime-specific design rule.}

\begin{itemize}\setlength{\itemsep}{2pt}\setlength{\parskip}{0pt}
\item @frison1992 found a single pre-treatment ANCOVA covariate can reduce power at low pre-post correlation.
\item The same degree-of-freedom cost applies in the run-in setting.
\item Within the intermediate band of Table \\ref{tab:anomaly-boundary}, designers who add run-in observations should plan at least two rather than exactly one.
\end{itemize}

\subsection{Correlation structure governs diminishing returns}\label{correlation-structure-governs-diminishing-returns}

\textbf{Under compound symmetry diminishing returns stem from the information matrix, whereas under AR(1) declining correlation accelerates them.}

\begin{itemize}\setlength{\itemsep}{2pt}\setlength{\parskip}{0pt}
\item The run-in change-score correlation is constant in $J_0$ under compound symmetry.
\item Diminishing returns for $J_0 \ge 2$ therefore arise from the information matrix structure.
\item Under AR(1) the correlation falls with temporal distance, yielding faster diminishing returns from distant run-in observations.
\end{itemize}

\section{\texorpdfstring{Optimal allocation for fixed \(J\)}{Optimal allocation for fixed J}}\label{optimal-allocation-for-fixed-j}

\textbf{Under a fixed total budget the optimum never allocates observations to the run-in phase.}

\begin{itemize}\setlength{\itemsep}{2pt}\setlength{\parskip}{0pt}
\item In every cell of Table \\ref{tab:optimal-grid}, $J_0^* = 0$; the budget splits between treatment and common close.
\item The run-in phase adds the pre-randomization slope parameter $\delta$ without contributing to the treatment contrast, so a run-in visit is dominated by a treatment or close visit whenever visits are fungible.
\item Paper 1's run-in gains are not contradicted: they arise when run-in visits are additional, appended to a fixed treatment phase, not reallocated from it.
\end{itemize}

\section{When is the common close phase beneficial?}\label{when-is-the-common-close-phase-beneficial}

\textbf{The common close phase aids estimation only indirectly, through the random slope, so the meaningful question is the fixed-budget trade-off.}

\begin{itemize}\setlength{\itemsep}{2pt}\setlength{\parskip}{0pt}
\item Common close observations inform the random slope $b_i$; they do not contribute to the treatment comparison directly.
\item Added to a fixed design, a common close observation reduced $\Var(\hat\gamma)$ in every configuration examined ($J_1 = 1, \ldots, 4$, $\sigma_b/\sigma$ down to $10^{-3}$), so no benefit threshold exists for free observations.
\item The design-relevant question is whether a budgeted visit is better spent on the common close than on the treatment phase.
\end{itemize}

\subsection{\texorpdfstring{Crossover as a function of the budget \(J\)}{Crossover as a function of the budget J}}\label{crossover-as-a-function-of-the-budget-j}

\textbf{The crossover ratio falls as the budget grows, and above it the close phase absorbs an increasing share of the budget.}

\begin{itemize}\setlength{\itemsep}{2pt}\setlength{\parskip}{0pt}
\item The crossover decreases from 0.81 at $J = 3$ to 0.30 at $J = 8$.
\item Above the crossover the full optimizer allocates a growing share to the close, reaching an even split $(0, 4, 4)$ at $J = 8$ at every examined ratio from 0.6 upward.
\item Even at $\sigma_b/\sigma = 0.32$ the optimum at $J = 8$ is $(0, 5, 3)$, so the close enters at low ratios once the budget is large.
\end{itemize}

\subsection{General time points}\label{general-time-points}

\textbf{Maximally separated time points minimize the slope-estimation variance, consistent with the classical extremes result.}

\begin{itemize}\setlength{\itemsep}{2pt}\setlength{\parskip}{0pt}
\item Within the raw-outcome model, the variance is minimized when $t_1 t_2$ is most negative.
\item This occurs when the two time points are maximally separated.
\item It is consistent with the known optimality of observations at the extremes of follow-up for slope estimation.
\end{itemize}

\subsection{Motivation}\label{motivation}

\textbf{The averaging logic that motivates run-in can be applied symmetrically at the end of treatment to reduce endpoint measurement error.}

\begin{itemize}\setlength{\itemsep}{2pt}\setlength{\parskip}{0pt}
\item Run-in averages pre-randomization observations to cut the measurement-error contribution to the baseline slope.
\item Replicated measures raise power with diminishing returns past 20-50 replications.
\item Closely spaced replicates of the final visit average out endpoint measurement error under the same treatment condition.
\end{itemize}

\textbf{Replicating an endpoint or baseline visit scales down its residual variance while leaving the random-slope contribution intact.}

\begin{itemize}\setlength{\itemsep}{2pt}\setlength{\parskip}{0pt}
\item Replicating the last visit $n_{\text{rep}}$ times over a short interval gives effective residual variance $\sigma^2 / n_{\text{rep}}$.
\item The random slope contribution $\sigma_b^2 t^2$ is essentially unchanged.
\item Replicating the baseline visit similarly reduces $\sigma^2$ in the change-score covariance.
\end{itemize}

\textbf{Replicated endpoints differ fundamentally from common close observations in timing and the information they provide.}

\begin{itemize}\setlength{\itemsep}{2pt}\setlength{\parskip}{0pt}
\item Common close observations occur at new time points after treatment stops.
\item They inform the off-treatment slope.
\item Replicated endpoints sit at the same time point as the last treatment visit and only reduce measurement error.
\end{itemize}

\subsection{Variance formula with replication}\label{variance-formula-with-replication}

\textbf{Replication enters only through modified entries of $R$, leaving the Woodbury computation otherwise unchanged.}

\begin{itemize}\setlength{\itemsep}{2pt}\setlength{\parskip}{0pt}
\item The diagonal entry for a replicated visit uses $\sigma^2 / n_{\text{rep}}$ in place of $\sigma^2$.
\item Off-diagonal entries reflecting shared baseline error use $\sigma^2 / n_{\text{rep,pre}}$.
\item The remaining Woodbury computation proceeds identically.
\end{itemize}

\textbf{Endpoint replication helps most when measurement error dominates and little when slope variability does.}

\begin{itemize}\setlength{\itemsep}{2pt}\setlength{\parskip}{0pt}
\item The benefit is largest at small $\sigma_b / \sigma$.
\item At large $\sigma_b / \sigma > 2$ the gain is modest.
\item Variance is then dominated by $\sigma_b^2 t^2$, which replication cannot reduce.
\end{itemize}

\subsection{Replication vs additional run-in observations}\label{replication-vs-additional-run-in-observations}

\textbf{A fixed budget of extra measurements can be allocated either to run-in observations or to endpoint replicates.}

\begin{itemize}\setlength{\itemsep}{2pt}\setlength{\parskip}{0pt}
\item Run-in observations sit at distinct time points separated by $t$.
\item Endpoint replicates are closely spaced at one time point.
\item Replicates provide measurement-error reduction only, not slope information.
\end{itemize}

\textbf{Whether run-in or endpoint replication wins depends on the signal-to-noise ratio and on the number of extra measurements.}

\begin{itemize}\setlength{\itemsep}{2pt}\setlength{\parskip}{0pt}
\item At $\sigma_b / \sigma = 1$, a single extra measurement is better spent on a replicate (variance 2.73 versus 2.82); run-in wins from the second extra measurement onward.
\item Run-in observations supply information about the individual slope $b_i$, not only measurement-error reduction, but a lone run-in observation pays the single-run-in penalty of Section 3.
\item At $\sigma_b / \sigma = 0.5$, replication won at every budget examined (one to four extra measurements) because measurement error dominates the variance.
\end{itemize}

\subsection{Combined replication at both ends}\label{combined-replication-at-both-ends}

\textbf{Replicating both terminal visits yields symmetric measurement-error reduction.}

\begin{itemize}\setlength{\itemsep}{2pt}\setlength{\parskip}{0pt}
\item Baseline replication reduces $\sigma^2$ at the start of follow-up.
\item Endpoint replication reduces $\sigma^2$ at the end of treatment.
\item Together they lower measurement error symmetrically at both ends.
\end{itemize}

\subsection{Pattern-mixture approach}\label{pattern-mixture-approach}

\textbf{Dropout is accommodated by a pattern-mixture model that combines stratum-specific sample sizes.}

\begin{itemize}\setlength{\itemsep}{2pt}\setlength{\parskip}{0pt}
\item The approach follows @frost2008 Section 2.3 and related pattern-mixture work.
\item Participants are stratified by their observed visit pattern.
\item The overall requirement is a weighted average across strata.
\end{itemize}

\textbf{Strata are truncated designs under monotone dropout, and the formula is a conservative approximation that exact treatments improve upon.}

\begin{itemize}\setlength{\itemsep}{2pt}\setlength{\parskip}{0pt}
\item Here $p_s$ is the proportion with pattern $s$ and $N_s$ the size needed if all had pattern $s$.
\item Under monotone dropout, the pattern after post-baseline visit $j$ is the truncated design retaining the first $j$ post-baseline observations; participants with no post-baseline visit contribute no slope-contrast information ($N_s = \infty$).
\item @hu2021 derive exact formulas under monotone dropout, more rigorous and less conservative when dropout is informative about the random slope.
\end{itemize}

\subsection{AD neuroimaging trial}\label{ad-neuroimaging-trial}

\textbf{The AD example evaluates optimal allocation under MIRIAD parameters across several budgets and dropout scenarios.}

\begin{itemize}\setlength{\itemsep}{2pt}\setlength{\parskip}{0pt}
\item Parameters are $\sigma^2 = 0.3363$, $\sigma_b^2 = 0.9745$, and $t = 1$ year.
\item Allocation is compared for $J = 4, 5, 6$ total observations.
\item Both with-dropout and without-dropout cases are considered.
\end{itemize}

\section{Discussion}\label{discussion}

\textbf{Under a fixed visit budget the optimum places nothing in the run-in phase, which reconciles rather than contradicts the run-in literature.}

\begin{itemize}\setlength{\itemsep}{2pt}\setlength{\parskip}{0pt}
\item In every configuration examined the optimizer sets $J_0^* = 0$ and splits the budget between treatment and common close.
\item Run-in visits pay for the pre-randomization slope parameter $\delta$ while informing the treatment contrast only indirectly, so a fungible visit is always better spent elsewhere.
\item Paper 1's efficiency gains apply when run-in visits are additional; the two framings answer different design questions.
\end{itemize}

\textbf{The signal-to-noise ratio and the budget jointly determine how much of the budget the common close should absorb.}

\begin{itemize}\setlength{\itemsep}{2pt}\setlength{\parskip}{0pt}
\item Reallocating one visit from treatment to common close pays once $\sigma_b/\sigma$ exceeds a crossover that falls from 0.81 at $J = 3$ to 0.30 at $J = 8$.
\item Above the crossover the close phase absorbs a growing share, reaching an even split at $J = 8$ in the settings examined.
\item Designers should estimate the slope-to-noise ratio from pilot data before fixing the phase split.
\end{itemize}

\textbf{Two refinements, unequal spacing and endpoint replication, offer further efficiency without changing the phase counts.}

\begin{itemize}\setlength{\itemsep}{2pt}\setlength{\parskip}{0pt}
\item In the raw-outcome special case of Section 6, spacing two observations toward the extremes of the window minimizes the slope-estimation variance; the general change-score case remains open.
\item Replicating the endpoint or baseline visit reduces measurement error and is most valuable when measurement error dominates.
\item At $\sigma_b/\sigma = 1$, a single extra measurement favors an endpoint replicate; run-in observations at new time points win from the second extra measurement onward.
\end{itemize}

\textbf{The allocation rules are derived under idealizing assumptions whose relaxation is the natural next step.}

\begin{itemize}\setlength{\itemsep}{2pt}\setlength{\parskip}{0pt}
\item The optima assume equally spaced observations, a linear trajectory, and compound-symmetric residuals.
\item Dropout is handled by a pattern-mixture approximation over monotone truncation strata, which the exact treatment of @hu2021 would refine.
\item Under AR(1) residuals the diminishing returns to distant run-in observations are steeper, shifting the optimum toward fewer, closer run-in observations.
\end{itemize}

\section{Conflict of interest}\label{conflict-of-interest}

\textbf{The author reports no conflicts of interest.}

\begin{itemize}\setlength{\itemsep}{2pt}\setlength{\parskip}{0pt}
\item No financial conflicts are present.
\item No non-financial conflicts are present.
\item No competing interests bear on this work.
\end{itemize}

\section{Funding}\label{funding}

\textbf{The work received no specific funding.}

\begin{itemize}\setlength{\itemsep}{2pt}\setlength{\parskip}{0pt}
\item No external grant supported the study.
\item No internal grant supported the study.
\item No sponsor influenced the analysis.
\end{itemize}

\section{Data availability statement}\label{data-availability-statement}

\textbf{The implementing package, tests, and shared bibliography are openly available in the project repository.}

\begin{itemize}\setlength{\itemsep}{2pt}\setlength{\parskip}{0pt}
\item Allocation routines reside in allocation.R and replicated.R; design.R builds the design matrices they rely on.
\item Validation tests and the four-paper bibliography are included.
\item Specific paths appear in the Reproducibility section below.
\end{itemize}

\section{Reproducibility}\label{reproducibility}

\textbf{The manuscript build relies on a small, fixed set of R packages.}

\begin{itemize}\setlength{\itemsep}{2pt}\setlength{\parskip}{0pt}
\item Rendering used R 4.4.x.
\item The in-package runinpower source and tinytest provide computation and testing.
\item rmarkdown and bookdown drive the manuscript build.
\end{itemize}

\textbf{All results are deterministic.}

\begin{itemize}\setlength{\itemsep}{2pt}\setlength{\parskip}{0pt}
\item Every table and figure is a deterministic evaluation of the variance formulas.
\item No Monte Carlo simulation appears in this manuscript.
\item No random seed is therefore required.
\end{itemize}

\textbf{Source code, companion manuscripts, and the shared bibliography are located at documented repository paths.}

\begin{itemize}\setlength{\itemsep}{2pt}\setlength{\parskip}{0pt}
\item The R package source is under R/.
\item Companion manuscripts (Papers 1, 2, 4) are under analysis/report/\{01,02,04\}-*/.
\item This manuscript source and the shared bibliography are under analysis/report/03-allocation/.
\end{itemize}

\end{document}
